
\uline{où} @@M0@@ \uline{et} @@M1@@
\uline{sont des immersions}, \uline{et} @@M2@@ \uline{un morphisme lisse}. \uline{Alors pour que} @@M3@@ \uline{soit une immersion régulière}, \uline{il faut et il suffit que} @@M4@@ \uline{soit une immersion régulière}.

i) Supposons @@M5@@ régulière. Posons @@M6@@. La donnée de @@M7@@ définit une section @@M8@@ de @@M9@@ au-dessus de @@M10@@, telle que @@M11@@:
@@M12@@

Or @@M13@@ est une immersion régulière, car c'est une section d'un morphisme lisse (VII 1.10); @@M14@@ est une immersion régulière car elle est déduite de @@M15@@ par un changement de base lisse, donc plat (VII 1.5). Donc @@M16@@ est une immersion régulière (VII 1.7).

ii) Supposons @@M17@@ régulière, et @@M18@@ étale. Alors @@M19@@ est une immersion ouverte (EGA IV 17.9.3), donc un isomorphisme de @@M20@@ sur @@M21@@; posons @@M22@@.

Soit @@M23@@ la restriction de @@M24@@ à @@M25@@, et soit @@M26@@ un ouvert de @@M27@@ tel que @@M28@@. On a alors : @@M29@@, donc @@M30@@ est une immersion régulière de @@M31@@ dans @@M32@@. Or le morphisme @@M33@@, étant lisse, est plat et de présentation finie, donc ouvert. C'est par conséquent un morphisme fidèlement plat et quasi-compact de @@M34@@ sur un voisinage ouvert de @@M35@@ dans @@M36@@. Il en résulte (VII 1.5) que @@M37@@ est une immersion régulière.

iii) Supposons @@M38@@ régulière, et que @@M39@@ soit un fibré vectoriel sur @@M40@@. Comme la propriété à prouver (@@M41@@ régulière) est locale sur @@M42@@, on peut supposer @@M43@@ affine, @@M44@@ fermée, et que @@M45@@. Posons @@M46@@,
